% Faithful translation baseline; Chinese source PDF pp. 25--30.
\section{Time Evolution and the Schr\"odinger Equation}
This chapter discusses the time evolution of states and physical quantities in quantum mechanics. Unlike other physical quantities, which are represented by operators, time in nonrelativistic quantum mechanics is only a parameter; it remains a classical quantity and is not an operator.

\subsection*{1) Time-evolution (translation) operator}
Initial state: $\ket{\alpha,t_0}$.

Evolution from $t_0$ to $t$:
\[\ket{\alpha,t}=\op U(t,t_0)\ket{\alpha,t_0}.\]
Expand $\ket{\alpha,t_0}$ and $\ket{\alpha,t}$ in eigenstates $\ket a$ of an operator $\op A$:
\[\ket{\alpha,t_0}=\sum_a\ket a\braket{a}{\alpha,t_0},\qquad
\ket{\alpha,t}=\sum_a\ket a\braket{a}{\alpha,t}.\]
Clearly, the probability of occupying a particular eigenstate is not necessarily conserved:
\[|\braket{a}{\alpha,t_0}|^2\ne|\braket{a}{\alpha,t}|^2.\]
However, the normalization condition
\[\braket{\alpha,t}{\alpha,t}=\braket{\alpha,t_0}{\alpha,t_0}=1\]
implies conservation of the total probability,
\[\sum_a|\braket{a}{\alpha,t}|^2=\sum_a|\braket{a}{\alpha,t_0}|^2=1,\]
and the time-evolution operator is unitary:
\[\op U^{\dagger}(t,t_0)\op U(t,t_0)=\mathbf1.\]
Require the composition law for time evolution:
\[\op U(t_2,t_0)=\op U(t_2,t_1)\op U(t_1,t_0),\qquad t_2>t_1>t_0.\]
An infinitesimal time-evolution operator can be written
\[\op U(t_0+\dd t,t_0)=\mathbf1-\ii\op\Omega\,\dd t.\]
Unitarity, $\op U^{\dagger}\op U=\mathbf1$, requires $\op\Omega$ to be Hermitian:
\[\op\Omega^{\dagger}=\op\Omega.\]
Since $\op U$ is dimensionless, $\op\Omega$ has dimensions of frequency. By the Planck--Einstein energy--frequency relation $E=\hbar\omega$, we may set
\[\op\Omega=\frac{\op H}{\hbar},\qquad
\boxed{\op U(t_0+\dd t,t_0)=\mathbf1-\frac{\ii}{\hbar}\op H\,\dd t.}\]

\subsection*{2) Equation of time evolution}
Using the composition law,
\[\op U(t+\dd t,t_0)=\op U(t+\dd t,t)\op U(t,t_0)
=\left(\mathbf1-\frac{\ii}{\hbar}\op H\dd t\right)\op U(t,t_0).\]
Hence
\[\frac{\op U(t+\dd t,t_0)-\op U(t,t_0)}{\dd t}
=-\frac{\ii}{\hbar}\op H\op U(t,t_0),\]
so that
\[\boxed{\ii\hbar\frac{\partial}{\partial t}\op U(t,t_0)
=\op H\op U(t,t_0).}\]
This is the fundamental equation of time evolution in quantum mechanics.

Multiplying the equation on the right by the initial state $\ket{\alpha,t_0}$ gives
\[\ii\hbar\frac{\partial}{\partial t}\op U(t,t_0)\ket{\alpha,t_0}
=\op H\op U(t,t_0)\ket{\alpha,t_0},\]
\[\boxed{\ii\hbar\frac{\partial}{\partial t}\ket{\alpha,t}=\op H\ket{\alpha,t}.}\]
This is the equation governing the time evolution of a state, the Schr\"odinger equation.

We now discuss solutions of this fundamental equation. If the Hamiltonian has no explicit time dependence, for example the magnetic interaction $\op H\sim\op J_z B$ in a constant magnetic field directed along $z$,
\[\boxed{\op U(t,t_0)=\exp\left[-\frac{\ii}{\hbar}\op H(t-t_0)\right].}\]
If the Hamiltonian depends on time but the Hamiltonians at different times commute, for example $\op H(t)\sim\op J_z B(t)$ when the magnitude but not the direction of the magnetic field changes,
\[\op U(t,t_0)=\exp\left[-\frac{\ii}{\hbar}\int_{t_0}^t\op H(t')\dd t'\right].\]
If the Hamiltonian depends on time and the Hamiltonians at different times do not commute, for example $\op H(t)\sim\vect{J}\cdot\vect{B}(t)$ when both the magnitude and direction of the field vary,
\[\begin{split}
\op U(t,t_0)=\mathbf1+\sum_{n=1}^{\infty}\left(-\frac{\ii}{\hbar}\right)^n
&\int_{t_0}^{t}\dd t_1\int_{t_0}^{t_1}\dd t_2\cdots
\int_{t_0}^{t_{n-1}}\dd t_n\\[-2pt]
&\qquad\times \op H(t_1)\op H(t_2)\cdots\op H(t_n).
\end{split}\]
This is called the Dyson series.

Below we mainly consider Hamiltonians that do not depend explicitly on time.

\subsection*{3) Stationary states}
Suppose initially the system is in an eigenstate of $\op H$:
\[\ket{\alpha,t_0}=\ket E,\qquad\op H\ket E=E\ket E.\]
Then
\[\ket{\alpha,t}=e^{-\ii\op H(t-t_0)/\hbar}\ket{\alpha,t_0}
=e^{-\ii E(t-t_0)/\hbar}\ket{\alpha,t_0}.\]
The state changes only by a phase. The expectation value of any physical quantity $\op A$ is
\[\mel{\alpha,t}{\op A}{\alpha,t}
=\mel{\alpha,t_0}{e^{\ii E(t-t_0)/\hbar}\op A e^{-\ii E(t-t_0)/\hbar}}{\alpha,t_0}
=\mel{\alpha,t_0}{\op A}{\alpha,t_0}\]
and does not change with time. Therefore, energy eigenstates are called stationary states.

If the system is initially not in a stationary state, its amplitude for occupying the stationary state $\ket E$ at time $t$ is
\[\braket{E}{\alpha,t}=\bra E e^{-\ii\op H(t-t_0)/\hbar}\ket{\alpha,t_0}
=e^{-\ii E(t-t_0)/\hbar}\braket{E}{\alpha,t_0},\]
and its probability satisfies
\[|\braket{E}{\alpha,t}|^2=|\braket{E}{\alpha,t_0}|^2.\]
The probability of occupying any stationary state is the same at all times as it was initially. Thus, if the system initially occupies a stationary state, it remains in that state; if its initial state is not stationary, it remains nonstationary.

\subsection*{4) Electron-spin precession}
Consider the interaction of an electron's spin magnetic moment with an external magnetic field directed along $z$:
\[\op H=-\left(\frac{e}{mc}\right)\vect{\op s}\cdot\vect B
=-\frac{eB}{mc}\op s_z=\omega\op s_z,
\qquad\omega=-\frac{eB}{mc}>0.\]
The time-evolution operator is
\[\op U(t,0)=e^{-\ii\op Ht/\hbar}=e^{-\ii\omega\op s_z t/\hbar}.\]
The operators $\op s_z$ and $\op H$ have common eigenstates $\ket{s_z^{\pm}}$:
\[\op s_z\ket{s_z^{\pm}}=\pm\frac{\hbar}{2}\ket{s_z^{\pm}},\qquad
\op H\ket{s_z^{\pm}}=\pm\frac{\hbar\omega}{2}\ket{s_z^{\pm}}.\]
Suppose the system initially occupies the positive eigenstate of $\op s_x$:
\[\ket{\alpha,0}=\ket{s_x^+}=\frac1{\sqrt2}\ket{s_z^+}+\frac1{\sqrt2}\ket{s_z^-}.\]
At an arbitrary time,
\[\ket{\alpha,t}=\op U(t,0)\ket{\alpha,0}.\]
The probability of finding the system in $\ket{s_x^+}$ is
\[\begin{split}
|\braket{s_x^+}{\alpha,t}|^2
&=\left|\left(\frac{\bra{s_z^+}+\bra{s_z^-}}{\sqrt2}\right)
 e^{-\ii\omega\op s_z t/\hbar}
 \left(\frac{\ket{s_z^+}+\ket{s_z^-}}{\sqrt2}\right)\right|^2\\
&=\left|\frac{e^{-\ii\omega t/2}+e^{\ii\omega t/2}}{2}\right|^2
=\cos^2\frac{\omega t}{2}.
\end{split}\]
By conservation of probability,
\[|\braket{s_x^-}{\alpha,t}|^2=1-|\braket{s_x^+}{\alpha,t}|^2
=\sin^2\frac{\omega t}{2}.\]
The expectation value at an arbitrary time is
\[\begin{split}
\langle s_x\rangle
&=\mel{\alpha,t}{\op s_x}{\alpha,t}
=\sum_{s_x}s_x|\braket{s_x}{\alpha,t}|^2\\
&=\frac{\hbar}{2}\left(\cos^2\frac{\omega t}{2}-\sin^2\frac{\omega t}{2}\right)
=\frac{\hbar}{2}\cos\omega t.
\end{split}\]
Similarly,
\[\langle s_y\rangle=\frac{\hbar}{2}\sin\omega t,\qquad\langle s_z\rangle=0.\]
\begin{center}
\begin{tikzpicture}[scale=1.05,>=Latex]
\draw[->](-1.4,0)--(1.45,0)node[right]{$y$};
\draw[->](0,-.85)--(0,2.1)node[above]{$z$};
\draw[->](0,0)--(-1.15,-.85)node[below]{$x$};
\draw (0,1.25) ellipse (1.0 and .32);
\draw[dashed] (0,0)--(-.76,1.46);
\draw[thick,->] (0,0)--(.72,1.49)node[right]{$\langle\vect s\rangle$};
\draw (.0,.35) arc[start angle=90,end angle=60,radius=.35];
\node at (.17,.62){$\alpha$};
\end{tikzpicture}
\end{center}
The expectation value $\langle\vect s\rangle$ precesses about the $z$ axis with time; the trajectory is a cone. This is the spin Larmor precession, with $\omega$ the Larmor frequency.

\subsection*{5) Energy--time uncertainty relation}
From
\[\langle(\Delta\op x)^2\rangle\langle(\Delta\op p)^2\rangle\ge\frac{\hbar^2}{4},\]
we obtain $\Delta x\,\Delta p\ge\hbar/2$, where
\[\Delta A=\sqrt{\langle(\Delta\op A)^2\rangle}.\]
By classical analogy, $x=x(t)$ and $p=p(E,p)$, one may suppose
\[\Delta t\,\Delta E\ge\frac\hbar2.\]
This is the energy--time uncertainty relation. However, unlike $x$, $p$, $E$, $\op x$, $\op p$, and $\op H$ in quantum mechanics, time $t$ remains a classical quantity. What does $\Delta t$ mean?

What are $\langle(\Delta\op H)^2\rangle$ and $\langle(\Delta t)^2\rangle$?

For an arbitrary observable $\op O$,
\[\langle O\rangle(t)=\mel{\alpha,t}{\op O}{\alpha,t}.\]
Its time derivative is
\[\frac{\dd}{\dd t}\langle O\rangle
=\left(\frac{\partial}{\partial t}\bra{\alpha,t}\right)\op O\ket{\alpha,t}
+\bra{\alpha,t}\frac{\partial\op O}{\partial t}\ket{\alpha,t}
+\bra{\alpha,t}\op O\left(\frac{\partial}{\partial t}\ket{\alpha,t}\right).\]
Since
\[\ii\hbar\frac{\partial}{\partial t}\ket{\alpha,t}=\op H\ket{\alpha,t},\qquad
-\ii\hbar\frac{\partial}{\partial t}\bra{\alpha,t}=\bra{\alpha,t}\op H,\]
we have
\[\frac{\dd}{\dd t}\langle O\rangle
=\frac1{\ii\hbar}\mel{\alpha,t}{[\op O,\op H]}{\alpha,t}
+\left\langle\frac{\partial\op O}{\partial t}\right\rangle.\]
If $\op O$ has no explicit time dependence,
\[\frac{\dd}{\dd t}\langle O\rangle
=\frac1{\ii\hbar}\langle[\op O,\op H]\rangle.\]
Consequently, the uncertainty relation between $\op O$ and $\op H$ gives
\[\langle(\Delta\op O)^2\rangle\langle(\Delta\op H)^2\rangle
\ge \left|\frac1{2\ii}\langle[\op O,\op H]\rangle\right|^2
=\frac{\hbar^2}{4}\left(\frac{\dd\langle O\rangle}{\dd t}\right)^2.\]
Therefore
\[\Delta O\,\Delta E\ge\frac\hbar2\left|\frac{\dd\langle O\rangle}{\dd t}\right|.\]
If we define
\[\Delta t=\Delta O\left|\frac{\dd\langle O\rangle}{\dd t}\right|^{-1},\]
then
\[\boxed{\Delta t\,\Delta E\ge\frac\hbar2.}\]
Thus $\Delta t$ is the time needed for the expectation value of a physical quantity $\op O$ to change by an amount equal to one standard deviation $\Delta O$. Clearly, both $\Delta t$ and $\Delta E$ are averages and uncertainties of physical quantities.

\textbf{Example.} In a stationary state $\Delta E=0$, and $\langle\op O\rangle$ is time independent:
\[\frac{\dd}{\dd t}\langle O\rangle=0,\qquad
\Delta t=\Delta O\left|\frac{\dd\langle O\rangle}{\dd t}\right|^{-1}\longrightarrow\infty.\]